\chapter*{ЗАКЛЮЧЕНИЕ}
\addcontentsline{toc}{chapter}{ЗАКЛЮЧЕНИЕ}

В результате выполнения выпускной квалификационной работы был разработан метод обеспечения согласованности состояний двух асинхронных игроков с помощью конечных автоматов.

В ходе выполнения работы были решены следующие задачи:
\begin{itemize}
	\item проведен анализ предметной области синхронизации игровых состояний и передачи данных в многопользовательских играх;
	\item рассмотрены существующие подходы к синхронизации состояний и способы передачи игровых данных;
	\item описаны типовые игровые данные, подлежащие хранению и передаче;
	\item определены основные причины рассинхронизации игровых состояний;
	\item разработан метод синхронизации состояний двух игроков с использованием конечных автоматов;
	\item определена формальная модель состояний системы, событий и переходов между состояниями;
	\item разработана структура конечного автомата, описывающего взаимодействие игроков;
	\item представлен алгоритм синхронизации состояний, включающий обработку входящих и исходящих событий, подтверждения доставки и повторную отправку сообщений;
	\item обоснован выбор Unity, языка C\# и UDP-сетевого обмена в качестве средств программной реализации;
	\item разработано программное обеспечение, реализующее предложенный метод синхронизации;
	\item разработано программное обеспечение для взаимодействия игроков на основе событий и конечных автоматов;
	\item проведено тестирование разработанного программного обеспечения;
	\item определены критерии оценки эффективности разработанного метода;
	\item разработана методика исследования эффективности метода при различных параметрах сетевого взаимодействия;
	\item выполнено сравнение разработанного метода с событийной синхронизацией без механизмов прикладной надежности.
\end{itemize}

Разработанный метод рассматривает синхронизацию не как простую передачу состояния между клиентом и сервером, 
а как задачу согласованного выполнения событий в распределенной асинхронной системе. Для этого логика взаимодействия игроков представлена в виде конечного автомата, 
а каждое значимое игровое действие обрабатывается как событие, допустимость которого проверяется относительно текущего состояния системы.

В программной реализации поверх UDP-сетевого обмена были реализованы механизмы прикладной надежности, такие как подтверждение обработки событий, 
хранение неподтвержденных событий, повторная отправка при отсутствии подтверждения, фильтрация дубликатов и контроль порядка сообщений. 
Эти механизмы позволяют снизить влияние потерь, дублирования и нарушения порядка доставки пакетов на корректность логического состояния игроков.
